\documentclass[a4paper,12pt]{article}
\usepackage[T2A]{fontenc}
\usepackage[utf8]{inputenc}
\usepackage[russian]{babel}
\usepackage{amsmath,amssymb,mathtools}
\renewcommand{\familydefault}{\sfdefault}
\usepackage{icomma}
\usepackage[a4paper,margin=18mm,top=17mm,bottom=19mm]{geometry}
\usepackage{microtype}
\usepackage{tikz}
\usetikzlibrary{calc,arrows.meta,positioning,shapes.geometric}
\usepackage{tabularx,booktabs,longtable,array}
\usepackage{enumitem}
\usepackage[hidelinks]{hyperref}
\usepackage{parskip}
\usepackage{fancyhdr}
\usepackage{setspace}
\usepackage{xcolor}

\definecolor{accent}{HTML}{2E6F73}
\definecolor{darkaccent}{HTML}{18484B}
\definecolor{textgray}{HTML}{2E3133}
\definecolor{softgray}{HTML}{EEF2F2}
\definecolor{warn}{HTML}{8A5A2B}

\color{textgray}
\onehalfspacing
\setlength{\parindent}{0pt}
\setlength{\parskip}{0.45em}
\setlist[itemize]{leftmargin=1.6em,itemsep=0.18em,topsep=0.2em}
\setlist[enumerate]{leftmargin=1.8em,itemsep=0.40em,topsep=0.25em}
\renewcommand{\arraystretch}{1.28}
\pagestyle{fancy}
\fancyhf{}
\fancyhead[L]{\small\color{darkaccent}Физика: движение по окружности}
\fancyhead[R]{\small\color{darkaccent}07.10.26}
\fancyfoot[C]{\small\thepage}
\setlength{\headheight}{15pt}

\newcommand{\sectionrule}{\par\vspace{0.25em}\noindent\color{accent}\rule{\linewidth}{0.55pt}\color{textgray}\par\vspace{0.25em}}
\newcommand{\key}[1]{\textbf{\color{darkaccent}#1}}
\newcommand{\formula}[1]{\[\boxed{#1}\]}
\newcommand{\Hz}{\,\text{Гц}}
\newcommand{\secnd}{\,\text{с}}
\newcommand{\metr}{\,\text{м}}
\newcommand{\ms}{\,\text{м/с}}
\newcommand{\mss}{\,\text{м/с}^2}
\newcommand{\rads}{\,\text{рад/с}}

\begin{document}

% ---------- TITLE ----------
\thispagestyle{empty}
\begin{center}
\vspace*{22mm}
{\Large\color{darkaccent}\textbf{Чек-лист}}\\[7mm]
{\Huge\bfseries Равномерное движение\\[2mm] по окружности}\\[8mm]
{\large Период, частота, скорости и центростремительное ускорение}\\[14mm]
\color{accent}\rule{0.62\linewidth}{0.8pt}\\[12mm]
{\Large Володя Хачатурян}\\[5mm]
{\large 7 октября 2026 г.}\\[22mm]
\begin{minipage}{0.80\linewidth}
\centering
\large
Главная цель: видеть связь между числом оборотов, временем, периодом и частотой, а затем уверенно переходить к линейной и угловой скорости и центростремительному ускорению.
\end{minipage}
\end{center}

\vfill
\begin{center}
{\small\color{darkaccent}Лёвин Артём Александрович эксклюзивно для Володи Хачатуряна}
\end{center}
\vspace{4mm}
\begin{tikzpicture}[remember picture,overlay]
\draw[accent,line width=1.2pt] ($(current page.south west)+(18mm,16mm)$)
  .. controls ($(current page.south west)+(55mm,28mm)$) and ($(current page.south east)+(-55mm,5mm)$)
  .. ($(current page.south east)+(-18mm,17mm)$);
\end{tikzpicture}
\clearpage

% ---------- CHECKLIST ----------
\section*{1. Что можно наблюдать напрямую}
\sectionrule

Пусть тело движется по окружности. За время $t$ оно совершило $n$ полных оборотов.

\begin{itemize}
  \item $n$ — число оборотов;
  \item $t$ — общее время наблюдения;
  \item $T$ — период: время одного полного оборота;
  \item $\nu$ — частота: число оборотов за единицу времени.
\end{itemize}

\formula{\nu=\frac{n}{t}}

Частота измеряется в герцах:
\[
1\Hz=1\,\text{с}^{-1}.
\]
Например, $\nu=3\Hz$ означает: тело совершает \key{3 оборота за 1 секунду}.

\formula{T=\frac{t}{n}}

Период измеряется в секундах. Например, $T=5\secnd$ означает: \key{один полный оборот длится 5 секунд}.

\section*{2. Связь периода и частоты}
\sectionrule

Формулы $\nu=n/t$ и $T=t/n$ взаимно обратны. Поэтому
\formula{\nu=\frac{1}{T},\qquad T=\frac{1}{\nu}}

\key{Проверка смысла:}
\begin{itemize}
  \item большой период $T$ означает медленные обороты и малую частоту $\nu$;
  \item большая частота $\nu$ означает много оборотов за секунду и малый период $T$.
\end{itemize}

\textbf{Короткие примеры.}
\[
T=2\secnd \Rightarrow \nu=\frac12=0{,}5\Hz,
\]
\[
\nu=10\Hz \Rightarrow T=\frac1{10}=0{,}1\secnd.
\]

\section*{3. Как быстро выражать величины из формул}
\sectionrule

На занятии использовался \key{диагональный перенос} — обычное свойство пропорции. Это особенно полезно в физике, где часто требуется выразить неизвестную величину.

Из
\[
\nu=\frac{n}{t}
\]
получаем
\formula{n=\nu t,\qquad t=\frac{n}{\nu}}

Из
\[
T=\frac{t}{n}
\]
получаем
\formula{t=nT,\qquad n=\frac{t}{T}}

\key{Практический совет:} если формула выглядит как дробь, сначала перепишите её как пропорцию, а затем аккуратно перенесите нужные величины по диагонали.

\clearpage
\section*{4. Линейная скорость}
\sectionrule

За один оборот тело проходит длину окружности
\[
L=2\pi R.
\]
За $n$ оборотов путь равен $s=2\pi Rn$. Тогда
\[
v=\frac{s}{t}=\frac{2\pi Rn}{t}.
\]
Так как $n/t=\nu$,
\formula{v=2\pi R\nu}

Используя $\nu=1/T$,
\formula{v=\frac{2\pi R}{T}}

Линейная скорость измеряется в $\text{м/с}$. Она показывает, какое расстояние вдоль окружности тело проходит за единицу времени.

\section*{5. Угловая скорость}
\sectionrule

Положение тела на окружности удобно описывать углом поворота $\varphi$. Один полный оборот соответствует
\[
360^\circ=2\pi\ \text{рад}.
\]

Угловая скорость — это скорость изменения угла:
\formula{\omega=\frac{\varphi}{t}}

За $n$ оборотов $\varphi=2\pi n$, поэтому
\[
\omega=\frac{2\pi n}{t}.
\]
Следовательно,
\formula{\omega=2\pi\nu=\frac{2\pi}{T}}

Единица угловой скорости — $\text{рад/с}$.

\section*{6. Связь линейной и угловой скорости}
\sectionrule

Сравним
\[
v=\frac{2\pi R}{T}
\qquad\text{и}\qquad
\omega=\frac{2\pi}{T}.
\]
Общий множитель $2\pi/T$ — это угловая скорость. Значит,
\formula{v=\omega R}

При одинаковой $\omega$ точка, расположенная дальше от центра, имеет большую линейную скорость.

\begin{center}
\begin{tikzpicture}[scale=1.08]
  \coordinate (O) at (0,0);
  \coordinate (A) at (3.1,1.4);
  \draw[accent,line width=1.1pt] (O) circle[radius=3.4];
  \draw[darkaccent,line width=1pt] (O)--(A) node[midway,below right] {$R$};
  \fill[darkaccent] (O) circle (2pt) node[below left] {$O$};
  \fill[accent] (A) circle (2.5pt) node[above right] {$A$};
  \draw[-{Latex[length=3mm]},accent,line width=1.2pt] (A) -- ++(-0.75,1.65) node[above] {$\vec v$};
  \draw[-{Latex[length=3mm]},warn,line width=1.2pt] (A) -- ($(A)!0.55!(O)$) node[midway,below right] {$\vec a_c$};
  \draw[-{Latex[length=3mm]},darkaccent,line width=1pt] (0.75,0) arc[start angle=0,end angle=27,radius=0.75];
  \node[darkaccent] at (1.02,0.26) {$\omega$};
\end{tikzpicture}
\end{center}

\key{Важно:} вектор линейной скорости направлен по касательной, а центростремительное ускорение — к центру окружности.

\clearpage
\section*{7. Центростремительное ускорение}
\sectionrule

При движении по окружности направление скорости непрерывно меняется. Поэтому даже при постоянном модуле скорости существует ускорение, направленное к центру окружности.

\formula{a_c=\frac{v^2}{R}}

Используя $v=\omega R$, можно получить ещё одну полезную форму:
\formula{a_c=\omega^2R}

Если вместо $v$ или $\omega$ подставить формулы через $T$ или $\nu$, можно выразить ускорение и через эти величины.

\key{Проверка единиц:}
\[
\frac{(\text{м/с})^2}{\text{м}}=\text{м/с}^2.
\]

\section*{8. Типичные ошибки}
\sectionrule
\begin{itemize}
  \item Путать частоту и период: $3\Hz$ — это 3 оборота в секунду; $T=3\secnd$ — один оборот за 3 секунды.
  \item Использовать минуты, километры или другие единицы без перевода в СИ, если ответ нужен в СИ.
  \item Путать $t$ и $T$: $t$ — общее время наблюдения, $T$ — время одного оборота.
  \item В формуле $v=2\pi R\nu$ забывать множитель $2\pi R$, то есть длину одного оборота.
  \item В угловой скорости использовать $360$ вместо $2\pi$ при расчётах в СИ. Для формул с $\omega$ угол берут в радианах.
  \item Считать центростремительное ускорение направленным по касательной. Оно всегда направлено \key{к центру окружности}.
\end{itemize}

\clearpage
% ---------- FLOWCHART ONLY PAGE ----------
\thispagestyle{plain}
\begin{center}
{\LARGE\bfseries\color{darkaccent}Алгоритм: задача на движение по окружности}
\end{center}
\vspace{8mm}

\tikzset{
 startstop/.style={ellipse,draw=accent,line width=1pt,minimum width=48mm,minimum height=12mm,align=center,font=\small},
 process/.style={rectangle,rounded corners=2mm,draw=accent,line width=1pt,minimum width=108mm,minimum height=13mm,text width=96mm,align=center,font=\small},
 decision/.style={diamond,aspect=2.0,draw=accent,line width=1pt,minimum width=56mm,minimum height=16mm,text width=45mm,align=center,font=\small},
 branch/.style={rectangle,rounded corners=2mm,draw=accent,line width=1pt,minimum width=66mm,minimum height=18mm,text width=58mm,align=center,font=\small},
 arr/.style={-{Latex[length=3mm]},line width=1pt,draw=darkaccent}
}
\begin{center}
\begin{tikzpicture}[node distance=9mm]
\node[startstop] (start) {Прочитайте условие и выпишите известные величины};
\node[process,below=of start] (si) {Переведите время, радиус и другие величины в СИ, если это требуется};
\node[decision,below=12mm of si] (base) {Известны $n,t$ или $T,\nu$?};
\node[branch,below left=16mm and 2mm of base] (nt) {Если известны $n,t$: найдите $\nu=n/t$ и при необходимости $T=1/\nu$.};
\node[branch,below right=16mm and 2mm of base] (tnu) {Если известны $T$ или $\nu$: используйте $\nu=1/T$.};
\node[process,below=20mm of $(nt)!0.5!(tnu)$] (choose) {Выберите нужную формулу: $v=2\pi R\nu$, $\omega=2\pi\nu$, $v=\omega R$, $a_c=v^2/R$.};
\node[process,below=of choose] (calc) {Подставьте числа, вычислите и проверьте единицы измерения};
\node[startstop,below=of calc] (finish) {Сопоставьте ответ с физическим смыслом};

\draw[arr] (start) -- (si);
\draw[arr] (si) -- (base);
\draw[arr] (base.west) -| node[pos=0.25,above,font=\small]{$n,t$} (nt.north);
\draw[arr] (base.east) -| node[pos=0.25,above,font=\small]{$T,\nu$} (tnu.north);
\draw[arr] (nt.south) |- (choose.west);
\draw[arr] (tnu.south) |- (choose.east);
\draw[arr] (choose) -- (calc);
\draw[arr] (calc) -- (finish);
\end{tikzpicture}
\end{center}
\clearpage

% ---------- TRAINING ----------
\section*{Тренировочный блок — Путь А}
\sectionrule
\textbf{10 заданий.} В решениях записывайте исходную формулу, преобразование при необходимости, подстановку и единицы.

\begin{enumerate}[label=\textbf{\arabic*.}]
  \item Тело совершило $60$ оборотов за $30\secnd$. Найдите частоту и период вращения.

  \item Период вращения равен $0{,}25\secnd$. Найдите частоту.

  \item Частота вращения равна $5\Hz$. Сколько оборотов совершит тело за $12\secnd$?

  \item Тело совершает $180$ оборотов за $1{,}5$ минуты по окружности радиуса $0{,}20\metr$. Найдите частоту и линейную скорость.

  \item Радиус окружности равен $0{,}30\metr$, период вращения $0{,}40\secnd$. Найдите линейную скорость.

  \item Частота вращения равна $3\Hz$. Найдите угловую скорость.

  \item Период вращения равен $0{,}50\secnd$. Найдите угловую скорость.

  \item Точка движется по окружности радиуса $0{,}40\metr$ с угловой скоростью $5\rads$. Найдите линейную скорость.

  \item Тело движется по окружности радиуса $2\metr$ со скоростью $8\ms$. Найдите центростремительное ускорение.

  \item За $50\secnd$ тело совершило $150$ оборотов по окружности радиуса $0{,}25\metr$. Найдите $\nu$, $T$, $\omega$, $v$ и $a_c$.
\end{enumerate}

\clearpage
% ---------- ANSWERS ----------
\section*{Ответы}
\sectionrule
\begin{longtable}{@{}>{\bfseries}p{0.10\linewidth}p{0.84\linewidth}@{}}
\toprule
№ & Ответ \\
\midrule
\endfirsthead
\toprule
№ & Ответ \\
\midrule
\endhead
1 & $\nu=2\Hz$, $T=0{,}5\secnd$. \\
2 & $\nu=4\Hz$. \\
3 & $n=60$ оборотов. \\
4 & $\nu=2\Hz$, $v=0{,}8\pi\ms\approx2{,}51\ms$. \\
5 & $v=1{,}5\pi\ms\approx4{,}71\ms$. \\
6 & $\omega=6\pi\rads\approx18{,}85\rads$. \\
7 & $\omega=4\pi\rads\approx12{,}57\rads$. \\
8 & $v=2\ms$. \\
9 & $a_c=32\mss$. \\
10 & $\nu=3\Hz$, $T=\dfrac13\secnd\approx0{,}33\secnd$, $\omega=6\pi\rads\approx18{,}85\rads$, $v=1{,}5\pi\ms\approx4{,}71\ms$, $a_c=9\pi^2\mss\approx88{,}83\mss$. \\
\bottomrule
\end{longtable}

\vfill
\begin{center}
\small\color{darkaccent}
Лёвин Артём Александрович эксклюзивно для Володи Хачатуряна
\end{center}

\end{document}
